\section{Problém splnitelnosti}\label{sec:general_sat}

Začneme zdánlivě možná jednoduchým problémem,~který však na poli informatiky hraje dosti důležitou roli. Předpokládejme,~že máme zadanou nějakou logickou formuli $\varphi$ o~libovolném počtu logických proměnných,~označme např. $x_1,x_2,\dots,x_n$. Naší otázkou je,~jestli je možné hodnoty proměnných $x_1,x_2,\dots,x_n$ nastavit tak (tj. dosadit hodnoty 0~a~1),~aby byla formule splněna\footnote{Může se na první pohled zdát,~že se jedná o~dosti teoretickou záležitost,~nicméně později uvidíme,~že mnoho jiných problémů má se \ensuremath{\SAT} silnou spojitost.},~tj. $\varphi(\dots)=1$. Takovému ohodnocení budeme říkat \emph{splňující}.

\problem{Splnitelnost obecné logické formule}{Logická formule $\varphi$.}{1,~pokud je $\varphi$ splnitelná,~jinak 0.}

Nejdříve si zopakujeme logické operace a~spojky. Mezi ně řadíme $\neg,~\land,~\lor,~\implies,~\iff$.
\begin{itemize}
    \item \textbf{Negace $\neg x$.} \emph{Převrací logickou hodnotu $x$.}
    \item \textbf{Konjunkce $a\land b$.} \emph{Pravdivá,~jsou-li obě hodnoty operandů $a,b$ pravdivé.}
    \item \textbf{Disjunkce $a\lor b$.} \emph{Pravdivá,~je-li alespoň jedna z~hodnot operandů $a,b$ pravdivá.}
    \item \textbf{Implikace $a\implies b$.} \emph{Je-li předpoklad $a$ splněn,~je splněn i~závěr $b$. Tj. implikace je nepravdivá,~pokud platí $a$ a~neplatí $b$.}
    \item \textbf{Ekvivalence $a\iff b$.} \emph{$a$ platí právě tehdy,~když platí $b$. Tj. ekvivalence je pravdivá,~jsou-li hodnoty operandů $a,b$ současně 0~nebo 1.}
\end{itemize}
\begin{table}[ht]\label{table:logicke_operace}
    \centering
    \begin{tabular}{|c|c|c|c|c|c|c|}
    \hline
    $a$ & $b$ & $\neg a$ & $a\land b$ & $a\lor b$ & $a\implies b$ & $a\iff b$ \\ \hline
    0   & 0   & 1        & 0          & 0         & 1             & 1         \\ \hline
    0   & 1   & 1        & 0          & 1         & 1             & 0         \\ \hline
    1   & 0   & 0        & 0          & 1         & 0             & 0         \\ \hline
    1   & 1   & 0        & 1          & 1         & 1             & 1         \\ \hline
    \end{tabular}
\end{table}
Zde konstatujme,~že ve všech příkladech bude mít negace $\neg$ vždy \emph{nejvyšší prioritu} oproti ostatním operacím. Prioritu ostatních operací budeme stanovovat pomocí závorek. 
\begin{example}\label{ex:sat_formule}
    \begin{itemize}
        \item $\varphi(x)=x\land\neg x$ \dots \textbf{ Není splnitelná},~pro $x=0$ i~$x=1$ je $\varphi=0$.
        \item $\psi(x)=x\lor\neg x$ \dots \textbf{ Je splnitelná},~a~to jak pro $x=0$,~tak $x=1$.
        \item $\chi(a,b,c)=\big((a\land b)\lor \neg c\big)\implies (\neg a\land\neg c)$ \dots \textbf{Je splnitelná},~např. pro $a=0$,~$b=0$ a~$c=0$.
    \end{itemize}
\end{example}
V~příkladu \ref{ex:sat_formule} výše se jedná o~poměrně jednoduché instance,~kdy jsme schopni vidět hned,~zda je formule splnitelná. Pokud bychom však uvážili nějakou složitější formuli,~problém se již značně zkomplikuje.
\importantbox{Naivní postup pro formuli o~$n$ proměnných prozkoumá až $2^n$ možných ohodnocení. Algoritmus zkoušející všechny možnosti je tak \emph{exponenciální}.}
Pro \ensuremath{\SAT} není dodnes známý žádný algoritmus,~který by jej uměl řešit pro libovolnou logickou formuli v~polynomiálním čase.

\subsection{Konjunktivní normální forma}\label{subsec:cnf}

Zkusíme se podívat na formule ve speciálním tvaru,~a~to tzv. \emph{konjunktivní normální formě},~neboli zkráceně \emph{CNF}.\footnote{Z~angl. \emph{conjunctive normal form}.}

Formule v~CNF se skládá z~\textbf{klauzulí},~které obsahují \textbf{literály}. Literálem je buď proměnná,~nebo její negace.
\begin{itemize}
    \item Klauzule jsou ve formuli odděleny \emph{logickou spojkou $\land$}
    \item a~v~každé klauzuli jsou literály odděleny \emph{logickou spojkou $\lor$}.
\end{itemize}
\[\varphi(\dots)=\underbrace{(\cdots\lor\cdots\lor\cdots)}_{\text{klauzule}}\land(\cdots\lor \overbrace{x}^{\text{literál}}\lor\cdots)\land(\cdots\lor\overbrace{\neg x}^{\text{literál}}\lor\cdots)\land(\cdots\lor \underbrace{y}_{\text{literál}}\lor\cdots)\land\cdots\]
\begin{example}\label{ex:cnf_formule}
    \begin{itemize}
        \item \(\varphi(x,y,z)=\neg x\land (y\lor z)\) \dots \textbf{Je v~CNF.} Formule $\varphi$ obsahuje klauzule $\neg x$ (klauzule o~jednom literálu) a~$y\lor z$.
        \item \(\psi(a,b)=a\land b\) \dots \textbf{Je v~CNF.}
        \item \(\chi(a,b,c,d)=a\land \big(b\lor (c\land d)\big)\) \dots \textbf{Není v~CNF},~protože výraz $c\land d$ není literál.
    \end{itemize}
    Poslední formuli lze však převést do CNF: $a\land (b\lor c)\land (b\lor d)$.
\end{example}
Z~příkladu \ref{ex:cnf_formule} se nabízí otázka,~zda je možné \emph{libovolnou formuli} převést do CNF. Existuje pro každou formuli ekvivalentní\footnote{Formule $\varphi$ a~$\psi$ jsou ekvivalentní,~pokud mají při každém ohodnocení společných proměnných stejnou pravdivostní hodnotu.} formule v~CNF? Ano.
\begin{theorem}
    Pro každou logickou formuli $\psi$ existuje ekvivalentní formule $\psi^\prime$ v~CNF.
\end{theorem}
Toto tvrzení lze,~pochopitelně,~dokázat,~ale zde se tím zabývat nebudeme a~zkrátka jej přijmeme jako fakt. Podstatná věc,~která z~toho plyne,~je,~že pokud je nějaká formule $\psi$ splnitelná,~pak je splnitelná i~jí ekvivalentní formule $\psi^\prime$ v~CNF a~naopak pokud $\psi$ není splnitelná,~není splnitelná ani $\psi^\prime$. Symbolicky
\[\psi\,\text{je splnitelná}\iff\psi^\prime\,\text{je splnitelná}.\]

\subsection{Převod do CNF pomocí tabulky}\label{subsec:cnf_prevod_tabulka}

Podívejme se na způsob,~jak lze převést libovolnou logickou formuli do CNF.

Máme obecnou formuli $\varphi$ s~proměnnými $x_1,x_2,\dots,x_n$,~pro kterou budeme konstruovat ekvivalentní formuli $\varphi^\prime$ v~CNF. Vycházíme z~tabulky pravdivostních hodnot,~přičemž pro každé \textbf{nepravdivé} ohodnocení vytvoříme \emph{klauzuli} s~$n$ literály,~kde obecně $i$-tý literál je
\begin{itemize}
    \item $x_i$,~pokud v~daném ohodnocení je $x_i=0$,
    \item jinak $\neg x_i$ (tj. když $x_i=1$).
\end{itemize}
Podívejme se na úvod na příklad \ref{ex:prevod_cnf_1} níže.
\begin{example}\label{ex:prevod_cnf_1}
    Máme formuli $\psi(x,y,z)=(x\implies y)\implies z$,~pro kterou chceme najít ekvivalentní formuli $\psi^\prime$ v~CNF. Nejprve si tedy sestavíme tabulku pravdivostních hodnot.
    \begin{table}[ht]
        \centering
        \begin{tabular}{|c|c|c|c|c|}
        \hline
        $x$ & $y$ & $z$ & $x\implies y$ & $(x\implies y)\implies z$ \\ \hline
        0   & 0   & 0   & 1             & 0                         \\ \hline
        0   & 0   & 1   & 1             & 1                         \\ \hline
        0   & 1   & 0   & 1             & 0                         \\ \hline
        0   & 1   & 1   & 1             & 1                         \\ \hline
        1   & 0   & 0   & 0             & 1                         \\ \hline
        1   & 0   & 1   & 0             & 1                         \\ \hline
        1   & 1   & 0   & 1             & 0                         \\ \hline
        1   & 1   & 1   & 1             & 1                         \\ \hline
        \end{tabular}
    \end{table}
    Z~tabulky můžeme vidět,~že formule $\psi$ je nepravdivá pro
    \begin{itemize}
        \item $x=0,y=0,z=0$,
        \item $x=0,y=1,z=0$
        \item a~$x=1,y=1,z=0$.
    \end{itemize}
    Tedy celkově sestavíme 3~klauzule:
    \begin{itemize}
        \item \((\overbrace{x}^{x=0}\lor \underbrace{y}_{y=0}\lor \overbrace{z}^{z=0})\),
        \item \((x\lor\underbrace{\neg y}_{y=1}\lor z)\),
        \item \((\neg x\lor\neg y\lor z)\).
    \end{itemize}
    Výsledná formule v~CNF bude mít tvar:
    \[\psi^\prime(x,y,z)=(x\lor y\lor z)\land (x\lor \neg y\lor z)\land (\neg x\lor \neg y\lor z).\]
    Čtenář se sám může přesvědčit,~že formule $\psi$ a~$\psi^\prime$ pro libovolné ohodnocení mají stejnou pravdivostní hodnotu.
\end{example}
Nabízí se otázka: \emph{\uv{Proč tento postup funguje?}} Pokud formule v~CNF není při určitém ohodnocení pravdivá,~alespoň jedna její klauzule není pravdivá. Klauzule obsahující všechny proměnné je nepravdivá právě při jediném ohodnocení: všechny její literály musí být nepravdivé. Pro každé ohodnocení,~při němž je původní formule nepravdivá,~tedy sestrojíme klauzuli,~která zakáže právě toto ohodnocení.
\begin{example}
    \(\chi(x_1,x_2,x_3,x_4)=(x_1\iff x_2)\lor \neg(x_3\iff x_4)\).
    \begin{table}[ht]
        \centering
        \begin{tabular}{|c|c|c|c|c|c|c|}
        \hline
        $x_1$ & $x_2$ & $x_3$ & $x_4$ & $x_1\iff x_2$ & $\neg(x_3\iff x_4)$ & $(x_1\iff x_2)\lor\neg(x_3\iff x_4)$ \\ \hline
        0     & 0     & 0     & 0     & 1             & 0                                                 & 1                                    \\ \hline
        0     & 0     & 0     & 1     & 1             & 1                                                 & 1                                    \\ \hline
        0     & 0     & 1     & 0     & 1             & 1                                                 & 1                                    \\ \hline
        0     & 0     & 1     & 1     & 1             & 0                                                 & 1                                    \\ \hline
        0     & 1     & 0     & 0     & 0             & 0                                                 & 0                                    \\ \hline
        0     & 1     & 0     & 1     & 0             & 1                                                 & 1                                    \\ \hline
        0     & 1     & 1     & 0     & 0             & 1                                                 & 1                                    \\ \hline
        0     & 1     & 1     & 1     & 0             & 0                                                 & 0                                    \\ \hline
        1     & 0     & 0     & 0     & 0             & 0                                                 & 0                                    \\ \hline
        1     & 0     & 0     & 1     & 0             & 1                                                 & 1                                    \\ \hline
        1     & 0     & 1     & 0     & 0             & 1                                                 & 1                                    \\ \hline
        1     & 0     & 1     & 1     & 0             & 0                                                 & 0                                    \\ \hline
        1     & 1     & 0     & 0     & 1             & 0                                                 & 1                                    \\ \hline
        1     & 1     & 0     & 1     & 1             & 1                                                 & 1                                    \\ \hline
        1     & 1     & 1     & 0     & 1             & 1                                                 & 1                                    \\ \hline
        1     & 1     & 1     & 1     & 1             & 0                                                 & 1                                    \\ \hline
        \end{tabular}
    \end{table}
    \[\chi^\prime(x_1,x_2,x_3,x_4)=(x_1\lor\neg x_2\lor x_3\lor x_4)\land(x_1\lor\neg x_2\lor\neg x_3\lor\neg x_4)\land(\neg x_1\lor x_2\lor x_3\lor x_4)\land(\neg x_1\lor x_2\lor\neg x_3\lor\neg x_4).\]
    \importantbox{\textbf{Problém:} Generování tabulky (resp. procházení všech možných ohodnocení) nás přivádí zpět k~původnímu problému,~a~to sice \emph{exponenciální časové složitosti}.}
    Tento postup nelze obecně zachránit pouhou chytřejší implementací: pokud trváme na ekvivalentní CNF bez nových proměnných,~může být výsledná formule exponenciálně delší. Například při roznásobení formule
    \[(x_1\land y_1)\lor(x_2\land y_2)\lor\dots\lor(x_n\land y_n)\]
    vzniká pro každou volbu jednoho literálu z~každé dvojice samostatná klauzule,~tedy celkem $2^n$ klauzulí.
\end{example}

\subsection{Tseitinova transformace}\label{subsec:cnf_tseitin}

Značně rychlejší převod do CNF navrhl sovětský matematik \emph{Grigorij S.~Cejtin}\footnote{Jeho příjmení se v~anglických textech obvykle přepisuje jako \emph{Tseitin}; odtud pochází ustálený název transformace.}. Předchozí postup trval na tom,~že ve výsledku ponecháme pouze původní proměnné. Tseitinova transformace zavede pomocné proměnné a~získá formuli,~která sice obecně není s~původní formulí ekvivalentní nad rozšířenou množinou proměnných,~ale je s~ní \emph{stejně splnitelná}.

Postup lze rozdělit do čtyř kroků.
\begin{enumerate}
    \item Formuli rozložíme na podformule odpovídající uzlům jejího syntaktického stromu.
    \item Pro každou složenou podformuli zavedeme novou proměnnou.
    \item Pro každou novou proměnnou přidáme klauzule,~které vynutí,~aby měla stejnou hodnotu jako příslušná podformule.
    \item Nakonec přidáme jednotkovou klauzuli tvořenou proměnnou zastupující celou původní formuli; tím požadujeme,~aby její hodnota byla $1$.
\end{enumerate}

Potřebné klauzule jsou krátké a~lze je zapsat přímo. Pro literály nebo pomocné proměnné $a,b$ platí
\begin{align}
y\iff\neg a
    &\quad\leadsto\quad(\neg y\lor\neg a)\land(y\lor a),\label{eq:tseitin-neg}\\
y\iff(a\land b)
    &\quad\leadsto\quad(\neg y\lor a)\land(\neg y\lor b)\land(y\lor\neg a\lor\neg b),\label{eq:tseitin-and}\\
y\iff(a\lor b)
    &\quad\leadsto\quad(y\lor\neg a)\land(y\lor\neg b)\land(\neg y\lor a\lor b).\label{eq:tseitin-or}
\end{align}
Implikaci nejprve přepíšeme pomocí $a\implies b\equiv\neg a\lor b$. Každý výskyt spojky tak přidá nejvýše tři klauzule konstantní délky.

Pojďme se podívat na konkrétní příklad.
\begin{example}
    Máme formuli $\alpha(p,q,r)=((p\lor q)\land r)\implies \neg p$. Chceme pro ni sestrojit Tseitinovou transformací novou formuli $\beta$.
    
    Formule $\alpha$ obsahuje tyto dílčí podformule:
    \begin{itemize}
        \item $\neg p$,
        \item $p\lor q$,
        \item $(p\lor q)\land r$,
        \item a~$((p\lor q)\land r)\implies \neg p$ (původní formule $\alpha$).
    \end{itemize}
    Pro ně zavedeme proměnné $y_1,y_2,y_3$ a~$y_4$. Implikaci nejprve přepíšeme jako disjunkci:
    \begin{itemize}
        \item $y_4\iff\neg p$,
        \item $y_3\iff(p\lor q)$,
        \item $y_2\iff(y_3\land r)$
        \item a~$y_1\iff(\neg y_2\lor y_4)$.
    \end{itemize}
    Pomocí vztahů \eqref{eq:tseitin-neg}--\eqref{eq:tseitin-or} dostaneme formuli v~CNF:
    \begin{align*}
        \beta={}&y_1\\
        &\land(y_1\lor y_2)\land(y_1\lor\neg y_4)
            \land(\neg y_1\lor\neg y_2\lor y_4)\\
        &\land(\neg y_2\lor y_3)\land(\neg y_2\lor r)
            \land(y_2\lor\neg y_3\lor\neg r)\\
        &\land(y_3\lor\neg p)\land(y_3\lor\neg q)
            \land(\neg y_3\lor p\lor q)\\
        &\land(\neg y_4\lor\neg p)\land(y_4\lor p).
    \end{align*}

    Zkusme ohodnocení $p=1,q=0$ a~$r=0$. Původní formule je pravdivá:
    \[\alpha(1,0,0)=((1\lor 0)\land 0)\implies\neg 1=0\implies 0=1.\]
    Pomocné proměnné musí nabývat hodnot
    \begin{itemize}
        \item $y_4=0$,~protože $\neg p=0$,
        \item $y_3=1$,~protože $p\lor q=1$,
        \item $y_2=0$,~protože $y_3\land r=0$,
        \item a~$y_1=1$,~protože $\neg y_2\lor y_4=1$.
    \end{itemize}
    Po dosazení jsou všechny klauzule formule $\beta$ splněny.
\end{example}

Je-li původní formule splnitelná,~nastavíme každou pomocnou proměnnou podle hodnoty její podformule a~splníme všechny nové klauzule. Je-li naopak výsledná formule splnitelná,~klauzule vynucují správné hodnoty všech pomocných proměnných a~jednotková klauzule pro kořen říká,~že původní formule musí být pravdivá. Obě formule jsou tedy stejně splnitelné.
\tipbox{Každá spojka původní formule přidá jednu pomocnou proměnnou a~nejvýše tři klauzule. Velikost výsledné formule i~doba konstrukce jsou proto lineární ve velikosti původní formule.}
